В последовательных методах задача состоит в том, чтобы по одному
наблюдению $x_o$ эффективно оценить апостериорное распределение
$p(\theta\mid x_o)$. Вместо генерации всех обучающих симуляций из априорного
распределения алгоритм на каждом раунде использует proposal
$\tilde p_r(\theta)$, сосредоточенный в области параметров, которая
представляется перспективной. Из него генерируются пары
$\theta_i\sim\tilde p_r(\theta)$, $x_i\sim p(x\mid\theta_i)$, после чего
условная модель $q_\phi(\theta\mid x)$ обучается на накопленных данных.
Такой подход позволяет направлять вычислительный бюджет на параметры,
которые согласуются с наблюдением, но требует учитывать, что proposal
может отличаться от исходного априорного распределения.

В SMC-ABC апостериор приближается облаком частиц: параметры симулируются,
а частицы с симулированными данными, недостаточно близкими к наблюдению,
отбрасываются. Между популяциями порог расстояния уменьшается, а
оставшиеся частицы переиспользуются для предложения новых параметров.
Метод не обучает условную плотность и поэтому подходит как
симуляционный baseline, но требует выбрать расстояние и пороги принятия,
и один из очевидных минусов состоит в том, что число симуляций может
быстро расти, если вероятность принятия мала.

\subsection{Смещение при сэмплировании}

При обычном обучении максимизацией правдоподобия модель аппроксимирует
условное распределение, соответствующее proposal, из которого были
получены обучающие параметры:
\[
\tilde p(\theta\mid x)
=\frac{\tilde p(\theta)p(x\mid\theta)}{\tilde p(x)}.
\]
Искомое распределение определяется исходным априорным распределением:
\[
p(\theta\mid x)
=\frac{p(\theta)p(x\mid\theta)}{p(x)}
=\tilde p(\theta\mid x)
\frac{p(\theta)}{\tilde p(\theta)}
\frac{\tilde p(x)}{p(x)}.
\]
При фиксированном $x$ последний множитель не зависит от $\theta$ и
поглощается нормировкой, однако отношение $p(\theta)/\tilde p(\theta)$
в общем случае зависит от параметра. Поэтому без коррекции обучение на
данных из произвольного proposal приводит к смещению. Коррекция может
быть включена в функцию потерь, применена к обученной модели либо
устранена выбором proposal с подходящей структурой.

\subsection{Коррекция с помощью плотностей: SNPE-A и SNPE-B}

В SNPE-A \cite{papamakarios2016fast} смещение устраняется после обучения: сначала оценивается
$\tilde p(\theta\mid x)$, а затем результат умножается на отношение
$p(\theta)/\tilde p(\theta)$. В классическом варианте используются
смеси гауссовских распределений, для которых такое преобразование можно
выполнить аналитически. Это делает поправку явной, но связывает метод
с конкретным семейством моделей и требует подходящего вида как для
proposal, так и для апостериорной аппроксимации.

SNPE-B \cite{lueckmann2017flexible} переносит эту поправку в функцию потерь, используя веса
$w(\theta)=p(\theta)/\tilde p(\theta)$:
\[
\mathcal L_B(\phi)=
-\mathbb E_{\theta\sim\tilde p(\theta),\,x\sim p(x\mid\theta)}
\left[
\frac{p(\theta)}{\tilde p(\theta)}
\log q_\phi(\theta\mid x)
\right].
\]
В отличие от SNPE-A, этот подход не требует гауссовской смеси, но
нуждается в вычислении плотностей prior и proposal. Если proposal
назначает малую плотность параметрам, которые имеют заметную вероятность
при prior, веса могут иметь высокую дисперсию и затруднять обучение.

\subsection{SNPE-C: атомарная апостериорная оценка}

SNPE-C, также называемый APT \cite{greenberg2019apt}, на практике работает подобно классификатору.
Для каждого симулированного наблюдения $x_k$ строится набор
кандидатов параметров $\theta_1,\ldots,\theta_K$, и модель должна
назначить наибольшую оценку параметру, который породил это наблюдение.
В атомарном варианте оценки кандидатов пропорциональны
$q_\phi(\theta_j\mid x_k)/p(\theta_j)$:
\[
\mathcal L_C(\phi)=
-\mathbb E\left[
\log
\frac{q_\phi(\theta_k\mid x_k)/p(\theta_k)}
{\displaystyle\sum_{j=1}^{K}
q_\phi(\theta_j\mid x_k)/p(\theta_j)}
\right].
\]
Нормировка по конечному набору превращает оценку в задачу выбора
параметра, согласующегося с наблюдением. В библиотеке \texttt{sbi} этот подход
не требует явно вычислять плотность proposal и соответствующие веса,
поэтому удобен для последовательного обучения.
Однако лосс требует значений $\log q_\phi(\theta_j\mid x_k)$
для кандидатов и значения плотности prior. Таким образом, SNPE-C
устраняет необходимость в плотности proposal, но не необходимость в
оценке плотности апостериорной модели.

\subsection{Усечение prior}

Усечение proposal также используется в TSNPE \cite{deistler2022tsnpe}.
В Truncated SFMPE после первого раунда мы сужаем prior по области
$\mathcal S_r$, построенной из posterior предыдущего раунда. Для каждой
координаты параметра по сэмплам находим интервал, в который попадает
заданная доля значений, и расширяем его на величину \texttt{padding}.
В следующем раунде генерируем параметры из исходного prior и оставляем
те, что попали в эту область и допустимую область параметров. Поэтому
новый proposal — это prior,
ограниченный на $\mathcal S_r$:
\[
\tilde p(\theta)
=p(\theta\mid\theta\in\mathcal S_r)
=\frac{p(\theta)\mathbf 1_{\mathcal S_r}(\theta)}
{P(\theta\in\mathcal S_r)}.
\]

Мы не стали строить область по квантилям $\log q_\phi$: для FMPE вычисление
лог-плотности требует дополнительного решения ОДУ и оценки дивергенции
поля скоростей, что дорого при большом числе сэмплов. Вместо этого границы
задаются непосредственно по сэмплам posterior. Такой способ проще, но
область зависит от качества предыдущей оценки posterior; если prior редко
попадает в неё, отбраковка может потребовать много симуляций.
